\section{Introducción}
\subsection{Descripción del problema}
En la FISEI, la homologación de asignaturas requiere revisar planes analíticos, créditos y contenidos. La revisión manual demanda tiempo y puede dificultar la explicación de una resolución. La decisión académica corresponde al docente y debe conservar su criterio.

\subsection{Descripción del proyecto}
SIHA-IA (Sistema Inteligente de Homologación Académica) es una aplicación web para registrar solicitudes y apoyar su análisis mediante lectura documental, reglas configuradas y herramientas de inteligencia artificial local. El resultado sirve de apoyo para el personal académico; la decisión final es humana.

\subsection{Alcance de los Sprints 1, 2 y 3}
El Sprint 1 estableció la base: cuentas y roles, reglas, catálogo académico, API Gateway, expedientes y publicación de solicitudes. El Sprint 2 incorporó la lectura de planes, confirmación de datos dudosos, análisis de créditos y contenidos, y consulta de precedentes. El Sprint 3 integra el certificado SIHA único emitido por Secretaría, aplica la regla de créditos y contenidos aclarada con el docente, registra la decisión de docentes y administrador, considera movilidad externa, presenta una tabla de resultados y automatiza el registro manual del resultado.

\subsection{Objetivos}
\textbf{Objetivo general.} Apoyar el trámite de homologación con un sistema que registre y analice solicitudes y facilite una resolución académica trazable.

\textbf{Objetivos específicos.}
\begin{itemize}
\item Registrar expedientes y documentos, y conservar los datos de la solicitud.
\item Analizar créditos y contenidos según las reglas acordadas.
\item Presentar los hallazgos a docentes y permitir la resolución prevista por el flujo institucional.
\item Mantener el historial de resultados y decisiones.
\end{itemize}

\subsection{Marco normativo}
El proyecto considera los instructivos CAU-P-011, para cambio de carrera, y CAU-P-122, para movilidad. SIHA-IA apoya la revisión y el registro del caso, pero no sustituye el trámite ni la autoridad académica.

\subsection{Términos utilizados}
\begin{tabularx}{\textwidth}{|p{0.25\textwidth}|X|}\hline
\rowcolor{grisTarjeta}\textbf{Término} & \textbf{Significado en este informe}\\\hline
Certificado SIHA & Documento único emitido por Secretaría que contiene la información de la solicitud.\\\hline
Sílabo & Documento con información académica de una asignatura.\\\hline
Créditos & Valor académico asignado a una materia.\\\hline
Contenidos & Temas y unidades que se revisan para comparar asignaturas.\\\hline
Agente & Componente que procesa una solicitud y coordina el análisis.\\\hline
JWT & Token firmado que permite identificar la sesión y sus permisos.\\\hline
Gateway & Punto de entrada que valida y dirige las peticiones a los servicios.\\\hline
Microservicio & Servicio independiente que cumple una función del sistema.\\\hline
RabbitMQ & Servicio de mensajería usado para intercambiar eventos.\\\hline
Ollama & Herramienta que ejecuta modelos de inteligencia artificial localmente.\\\hline
CUMPLE & Resultado que indica que se alcanza la condición evaluada.\\\hline
NO\_CUMPLE & Resultado que indica que no se alcanza la condición evaluada.\\\hline
CRITERIO\_DOCENTE & Caso que requiere valoración y decisión de docentes.\\\hline
Bloqueada por créditos & Solicitud que no pasa a votación por no alcanzar el mínimo de créditos.\\\hline
\end{tabularx}
\captionof{table}{Términos utilizados}
